% Historical submitted candidate; superseded by the factored checkpoint91 text.
% See big-exponential-compactness-review.md for the completed review and exact scope.
% DRAFT: independent review pending; not an active book input.
\subsection{Limits of normalized exponential solutions}

The fixed-reference compactness theorem also describes the approximation
used to construct the zero-exponent equation. It controls differences of
potentials uniformly on $\Omega$, even when the potentials themselves
are unbounded near its complement.

\begin{proposition}\label{prop:big-exponential-clusters}
Let $0\leq f\in L^p(\mathrm{d}V)$, $p>1$, be nonzero, and let
$u_\lambda\simeq V_\theta$ be the unique solution of
\[
 \mu_{u_\lambda}=\mathrm{e}^{\lambda u_\lambda}f\,\mathrm{d}V,
 \qquad \lambda>0.
\]
Put
\[
 v_\lambda=u_\lambda-\sup_Xu_\lambda,
 \qquad c_\lambda=\mathrm{e}^{\lambda\sup_Xu_\lambda}.
\]
For $0<\lambda\leq1$ there are constants $B,C,m>0$, independent of
$\lambda$, such that
\[
 V_\theta-B\leq v_\lambda\leq V_\theta,
 \qquad m\leq c_\lambda\leq C.
\]
The normalized family is relatively compact for uniform convergence of
differences on $\Omega$. As $\lambda\downarrow0$, every cluster limit
$v$ is normalized, has minimal singularities, and satisfies
\[
 \mu_v=cf\,\mathrm{d}V,
\]
where $c>0$ is the unique scalar in
\cref{thm:big-minimal-normalized}. Moreover,
\[
 c_\lambda\longrightarrow c,
 \qquad c_\lambda\mathrm{e}^{\lambda v_\lambda}f
          \longrightarrow cf\quad\text{in }L^p(\mathrm{d}V).
\]
Every sequence of parameters in $(0,1]$ has a subsequence for which the
potentials converge in $C^\beta$ on each compact subset of $\Omega$, for
all
\[
 0<\beta<\frac{2(p-1)}{np+3(p-1)}.
\]
No metric mass property is assumed.
\end{proposition}

\begin{proof}
The equation for $v_\lambda$ is
\[
 \mu_{v_\lambda}=c_\lambda\mathrm{e}^{\lambda v_\lambda}
                         f\,\mathrm{d}V.
\]
Fix a Gauduchon metric $\gamma$ and write
$\gamma^n=G\,\mathrm{d}V$. Since $v_\lambda\leq0$ and
$0<\lambda\leq1$, \eqref{eq:big-singular-gauduchon-trace} gives
\[
 c_\lambda^{1/n}
 \int_X\mathrm{e}^{v_\lambda/n}f^{1/n}G^{1-1/n}\,\mathrm{d}V
 \leq\int_X\theta\wedge\gamma^{n-1}.
\]
The integrals on the left have a positive uniform lower bound.
Otherwise quasi-plurisubharmonic compactness would give an almost
everywhere convergent subsequence of the normalized $v_\lambda$ with a
genuine limit, and dominated convergence would give a positive limit
integral, since $f$ is nonzero. Thus $c_\lambda\leq C$.

The lower-type estimate in the proof of
\eqref{eq:big-normalized-uniform-type} now applies for every
$0<\lambda\leq1$. More explicitly, uniform Skoda integrability gives
fixed $\varepsilon>0$ and $1<q<p$ with
$\|\mathrm{e}^{-\varepsilon v_\lambda}f\|_q$ uniformly bounded.
The reference-preserving subsolution lemma supplies
\[
 V_\theta-1\leq\eta_\lambda\leq V_\theta,
 \qquad \mu_{\eta_\lambda}\geq
 m\mathrm{e}^{-\varepsilon v_\lambda}f\,\mathrm{d}V
 \geq mf\,\mathrm{d}V
\]
for a fixed $m>0$. Scalar order, together with
$\mu_{v_\lambda}\leq c_\lambda f\,\mathrm{d}V$, gives
$c_\lambda\geq m$. Set
$D=\varepsilon^{-1}\max\{0,\log(C/m)\}$ and
$w_\lambda=\eta_\lambda-D$. Then
\[
 \mathrm{e}^{\varepsilon(w_\lambda-v_\lambda)}\mu_{v_\lambda}
 \leq m\mathrm{e}^{-\varepsilon v_\lambda}f\,\mathrm{d}V
 \leq\mu_{w_\lambda},
\]
since $\lambda v_\lambda\leq0$ and $\eta_\lambda\leq0$.
Exponential order gives $v_\lambda\geq w_\lambda$, proving the uniform
bounds with $B=1+D$.

The densities are bounded in $L^p$ by $C\|f\|_p$. Apply
\cref{cor:big-reference-uniform-convergence} with $\eta=V_\theta$.
Given any sequence $\lambda_j\downarrow0$, pass to a subsequence with
$v_{\lambda_j}\to v$ uniformly in differences on $\Omega$ and in
$L^1(X)$, and with $c_{\lambda_j}\to c_*\in[m,C]$.
The limit has the same reference bounds and is normalized. On the
complement of the pole set,
\[
 0\leq1-\mathrm{e}^{\lambda v_\lambda}
 \leq1-\mathrm{e}^{\lambda(V_\theta-B)}\longrightarrow0.
\]
Dominated convergence after multiplication by $f$ therefore gives
$\mathrm{e}^{\lambda v_\lambda}f\to f$ in $L^p$ for the whole family.
The compactness corollary identifies the limit equation as
$\mu_v=c_*f\,\mathrm{d}V$. Scalar uniqueness gives $c_*=c$.
Every subsequential scalar limit as $\lambda\downarrow0$ is therefore $c$, so
$c_\lambda\to c$. The stated density convergence follows.

Finally \cref{thm:big-minimal-holder} gives uniform local $C^\alpha$
bounds for every $\alpha$ below its endpoint. For
$0<\beta<\alpha$, uniform convergence and the elementary interpolation
inequality on smaller coordinate balls give
\[
 [v_{\lambda_j}-v]_{C^\beta}
 \leq C\|v_{\lambda_j}-v\|_\infty^{1-\beta/\alpha}
          [v_{\lambda_j}-v]_{C^\alpha}^{\beta/\alpha}
 \longrightarrow0.
\]
Use slightly larger compact neighborhoods for these estimates. The
uniformly convergent subsequence already supplied by compactness works
on every compact subset and for each sub-endpoint exponent; no further
choice depending on that exponent is needed.
\end{proof}

\begin{corollary}\label{cor:big-exponential-symmetries}
Suppose a group $G$ of biholomorphisms preserves $\theta$ and the measure
$f\,\mathrm{d}V$. Every $v_\lambda$ and every cluster limit as
$\lambda\downarrow0$ is $G$-invariant. If the normalized zero-exponent
equation has only one $G$-invariant solution of minimal singularities,
then the whole family $v_\lambda$ converges to it, uniformly in
differences on $\Omega$ and in each local $C^\beta$ topology from the
proposition.
\end{corollary}
\begin{proof}
For $g\in G$, pullback preserves $\theta$-plurisubharmonicity and
$V_\theta\circ g=V_\theta$, by the envelope definition. It also
preserves the positive-exponent equation. Its uniqueness gives
$u_\lambda\circ g=u_\lambda$, and normalization preserves this equality.
Pullback by a fixed smooth biholomorphism is continuous on $L^1(X)$,
since its Jacobian relative to the fixed smooth volume is bounded.
Every cluster limit is therefore invariant almost everywhere, hence in
canonical representatives. This argument does not require $g$ to
preserve the chosen barrier or $\Omega$.

Under the stated uniqueness within the invariant class, every cluster
limit is the same potential. Relative compactness then implies
convergence of the whole family in the uniform-difference metric:
a sequence staying a fixed positive distance from that potential would
have a convergent subsequence, a contradiction. The interpolation
argument above upgrades it to the asserted local $C^\beta$ convergence.
\end{proof}

The proposition gives compactness and identifies every cluster equation.
It does not by itself select a unique cluster potential. The additional
hypothesis in the corollary concerns only the invariant class; uniqueness
of all normalized solutions can fail, as in
\cref{thm:degenerate-hermitian-nonuniqueness}.
